
\subsection{Sérialisabilité}


\begin{frame}{Définition de la sérialisabilité}


\begin{defblock}{Définition (à connaître et -- surtout -- comprendre:)}
Une exécution concurrente $H$ de $n$ transactions $T_1, \cdots, T_n$
est \alert{sérialisable} si et seulement si
il existe toujours un ordonnancement $H’$ de $T_1, \cdots, T_n$ tel que le résultat 
de l’exécution de $H$ est le même que celui de l'exécution \alert{en série} de $H’$.
\end{defblock}

\vfill

\begin{blocgauche}
Bien comprendre: si j'ai deux transactions $T_1$ et $T_2$, leur exécution imbriquée
est sérialisable ssi équivalente à $T_1$ puis $T_2$, ou à $T_2$ puis $T_1$ (\alert{en série}).

\medskip

En d'autres termes, à une exécution \alert{en isolation complète}
\end{blocgauche}

\end{frame}


\begin{frame}{Caractérisation de la sérialisabilité\,: conflits}

La définition précédente est \alert{déclarative}\,: elle dit ce qu'est la sérialisabilité,
pas la manière de la vérifier en pratique.

\medskip

Nous avons besoin de caractérisation plus opérationnelle.
\vfill

\begin{defblock}{Définition (à connaître)}
Deux opérations  $p_i[x]$ et $q_j[y]$ sont 
\alert{en conflit} si $x=y$, $i \not= j$,  $p$ ou 
$q$ est une écriture.
\end{defblock}

\vfill
\begin{blocgauche}
En clair: deux transactions accèdent au même nuplet, et (au moins) une veut le modifier.
\end{blocgauche}

\end{frame}

\begin{frame}{Exemple}

Reprenons une nouvelle fois l'exemple des mises à jour perdues:
$$r_1(s)  r_1(c_1) r_2(s)  r_2(c_2)  w_2(s)  w_2(c_2) w_1(s) w_1(c_1)$$

\begin{redbullet}
  \item $r_1(s)$ et $w_2(s)$ sont en conflit\,;
  \item $r_2(s)$ et $w_1(s)$ sont en conflit\,;
  \item $w_2(s)$ et $w_1(s)$ sont en conflit.
\end{redbullet}

\vfill

 $r_1(s)$ et $r_2(s)$ \alert{ne sont  pas} en conflit (lectures)\,; pas de conflit sur
$c_1$ et $c_2$.

\vfill
\begin{blocgauche}
Qui des conflits ici?
$$r_1(c_1) r_1(c_2) r_2(s)  r_2(c_2)  w_2(s)   w_2(c_2) r_1(s)$$
\end{blocgauche}

\end{frame}


\begin{frame}{Relation entre les transactions}

\begin{defblock}{Définition (à connaître)}
$H$  exécution concurrente des transactions $T = \{T_1, T_2, \cdots, T_n\}$.

Il existe une relation $\lhd$ sur
cet ensemble, définie par\,:
$$T_i \lhd T_j \Leftrightarrow \exists p \in T_i, q \in T_j, p\ est\ en\
conflit\ avec \ q\ et\ p <_H q $$
\noindent
où $p <_H q$ indique que $p$ apparaît avant $q$ dans $H$.
\end{defblock}

\vfill
\begin{blocgauche}
Dans l'exemple précédent: on a  $T_1 \lhd T_2$,
ainsi que $T_2 \lhd T_1$.
\end{blocgauche}
\end{frame}


\begin{frame}{Condition de sérialisabilité}
La condition sur la sérialisabilité
s'exprime sur le graphe de la relation $(T, \lhd)$,
dit
\alert{graphe de sérialisation}.

\vfill


\begin{defblock}{Condition de sérialisabilité}
Soit $H$ une exécution concurrente d'un ensemble de transactions
$T$. Si le graphe de $(T, \lhd)$ est acyclique, $H$ est sérialisable.
\end{defblock}

\vfill

\begin{blocgauche}
Autrement dit, $(T, \lhd)$ est une relation d'ordre partiel (antisymétrie).
\end{blocgauche}
\end{frame}


\begin{frame}{Graphes de sérialisabilité}

Lesquels correspondent à une exécution sérialisable?

\vfill

\figSlideWithSize{graphe-serial}{10cm}

\vfill

\begin{blocgauche}
\alert{Important}: un algorithme de contrôle de concurrence doit vérifier qu'aucun cycle ne peut intervenir
dans ce graphe.
\end{blocgauche}
\end{frame}

%%%%%%%%%%%%%%%%%%%%%%%%%%%%%%%%%%%%%%%%%%%%%%%%%%%%%%%
\begin{frame}{À retenir}

Caractérisation ``sémantique'' de la sérialisabilité\,: équivalence avec une exécution en série, pour au moins
un ordonnancement des transactions impliquées

\vfill

Caractérisation ``syntaxique'' de la sérialisabilité\,: pas de cycle dans le graphe de sérialisation
construit sur les conflits entre opérations.

\vfill

\alert{Objectif d'un algorithme de contrôle}\: s'assurer que toutes les exécutions sont sérialisables

\begin{blocgauche}
\begin{redbullet}
  \item en surveillant le graphe de sérialisation et en rejetant une transaction si un cycle apparaît (variante ``optimiste'')
  \item en tentant de prévenir l'apparition de cycles (variante ``pessimiste'')
\end{redbullet}
\end{blocgauche}

\end{frame}
